\hypertarget{ej6}{}
\noindent \textbf{Ejercicio 6.} Asumiendo que el valor de verdad de $b$ y $c$ es verdadero, el de $a$ es falso y el de $x$ e $y$ es indefinido, indicar cuáles de los operadores deben ser operadores "luego" para que la expresión no se indefina nunca:
\begin{multicols}{2}
\begin{enumerate}[label=\alph*)]
    \item $(\neg x \vee b)$
    \item $((c \vee (y \wedge a)) \vee b)$
    \item $\neg(c \vee y)$
    \item $(\neg(c \vee y) \leftrightarrow (\neg c \wedge \neg y))$
    \item $((c \vee y) \wedge (a \vee b))$
    \item $(((c \vee y) \wedge (a \vee b)) \leftrightarrow (c \vee (y \wedge a) \vee b))$
    \item $(\neg c \wedge \neg y)$
\end{enumerate}
\end{multicols}

\vspace{0.2cm}
{\color{gray}\hrule}
\vspace{0.4cm}

\textbf{Resolución:}

\textbf{Repaso teórico:} En lógica trivaluada (y en la programación en general), los operadores "luego" emplean \textbf{evaluación de cortocircuito}. Al evaluar estrictamente de izquierda a derecha, se detienen y devuelven un valor apenas el resultado ya está garantizado, evitando evaluar términos conflictivos o indefinidos ($\bot$):
\begin{itemize}
    \item \textbf{Disyunción luego ($\vee_L$):} Si el lado izquierdo es $V$, devuelve $V$ sin mirar el lado derecho ($V \vee_L \bot \equiv V$).
    \item \textbf{Conjunción luego ($\wedge_L$):} Si el lado izquierdo es $F$, devuelve $F$ sin mirar el lado derecho ($F \wedge_L \bot \equiv F$).
\end{itemize}

Valores dados: $b=V$, $c=V$, $a=F$, $x=\bot$, $y=\bot$. Analizamos caso por caso:

\begin{enumerate}[label=\textbf{\alph*)}]

    \vspace{0.3cm}
    \item \textbf{$(\neg x \vee b)$} \\
    Reemplazamos: $(\neg \bot \vee V)$. Al evaluar de izquierda a derecha, lo primero que el programa intenta hacer es resolver $\neg x$. Como $x$ es indefinido, la expresión falla de inmediato en ese punto. \\
    \underline{\textbf{Respuesta:}} No se puede evitar que se indefina (falla siempre). Ni siquiera un operador "luego" lo salva, porque el término indefinido está a la izquierda.

    \vspace{0.4cm}
    \item \textbf{$((c \vee (y \wedge a)) \vee b)$} \\
    Reemplazamos: $((V \vee (\bot \wedge F)) \vee V)$. Evaluamos de izquierda a derecha. En el primer paréntesis leemos la $c$ que vale $V$. Si el $\vee$ que le sigue es "luego", hace cortocircuito y todo el bloque $((V \vee_L (\bot \wedge F))$ se transforma en $V$ sin tocar la $y$. \\
    \underline{\textbf{Respuesta:}} El $\vee$ entre la $c$ y $(y \wedge a)$ debe ser un operador "luego".

    \vspace{0.4cm}
    \item \textbf{$\neg(c \vee y)$} \\
    Reemplazamos: $\neg(V \vee \bot)$. Para evitar que se evalúe la $y$ indefinida, necesitamos que la disyunción dentro del paréntesis haga cortocircuito al leer la $V$ inicial. \\
    \underline{\textbf{Respuesta:}} El operador $\vee$ debe ser un operador "luego".

    \vspace{0.4cm}
    \item \textbf{$(\neg(c \vee y) \leftrightarrow (\neg c \wedge \neg y))$} \\
    Analizamos ambos lados del bicondicional:
    \begin{itemize}
        \item \textbf{Lado izquierdo:} $\neg(V \vee \bot)$. Igual que en el inciso c), necesitamos que sea $\vee_L$ para que evalúe a $\neg(V) \equiv F$.
        \item \textbf{Lado derecho:} $(\neg V \wedge \neg \bot) \equiv (F \wedge \neg \bot)$. Al leer el primer valor, obtenemos un $F$. Si la conjunción es "luego", hace cortocircuito y devuelve $F$ de inmediato, evitando el $\neg \bot$.
    \end{itemize}
    \underline{\textbf{Respuesta:}} El $\vee$ del lado izquierdo y el $\wedge$ del lado derecho deben ser operadores "luego".

    \vspace{0.4cm}
    \item \textbf{$((c \vee y) \wedge (a \vee b))$} \\
    Reemplazamos: $((V \vee \bot) \wedge (F \vee V))$. El único peligro de indefinición está en el primer paréntesis. Como arranca con $V$, un $\vee_L$ lo hace valer $V$ automáticamente por cortocircuito. El segundo paréntesis no tiene indefiniciones, así que no importa qué operadores tenga. \\
    \underline{\textbf{Respuesta:}} El $\vee$ entre $c$ e $y$ debe ser un operador "luego".

    \vspace{0.4cm}
    \item \textbf{$(((c \vee y) \wedge (a \vee b)) \leftrightarrow (c \vee (y \wedge a) \vee b))$} \\
    Analizamos por partes:
    \begin{itemize}
        \item \textbf{Lado izquierdo:} Igual que en el inciso e), el primer $\vee$ (entre $c$ e $y$) debe ser $\vee_L$. Todo el bloque izquierdo terminará valiendo $V$.
        \item \textbf{Lado derecho:} $(V \vee (\bot \wedge F) \vee V)$. Evaluando de izquierda a derecha, la proposición arranca con $V$. Si el primer $\vee$ (el que le sigue a la $c$) es un $\vee_L$, hace cortocircuito y omite evaluar todo el resto de la expresión, devolviendo $V$.
    \end{itemize}
    \underline{\textbf{Respuesta:}} El primer $\vee$ del lado izquierdo y el primer $\vee$ del lado derecho deben ser operadores "luego".

    \vspace{0.4cm}
    \item \textbf{$(\neg c \wedge \neg y)$} \\
    Reemplazamos: $(\neg V \wedge \neg \bot) \equiv (F \wedge \neg \bot)$. Evaluando de izquierda a derecha, ya sabemos que el primer valor es falso ($F$). Si usamos una conjunción de cortocircuito ($\wedge_L$), la expresión entera devuelve $F$ sin intentar evaluar el $\neg y$. \\
    \underline{\textbf{Respuesta:}} El operador $\wedge$ debe ser un operador "luego".

\end{enumerate}

\vspace{0.5cm}
\hrule
\vspace{0.5cm}